% frontmatter.tex
\begin{titlepage}
\begin{center}

{\large School of Medicine, MSc Biomedical Informatics, 2024 to 2026\par}
\vspace{1.0cm}

{\large School of Medicine\par}
{\large MSc Biomedical Informatics\par}
\vspace{1.0cm}

{\large Author: Nikolaos Chaikalis\par}
{\large Supervisor: Professor Eleni Kaldoudi\par}
{\large Supervising PhD Candidate: Nikolaos Portoakalidis\par}

\vspace{1.8cm}

{\LARGE Diploma Thesis\par}
\vspace{1.2cm}

{\Large
Design and Evaluation of a Retrieval Augmented Biomedical Platform for Thrombosis Risk Factor Extraction: Continuous PubMed Scheduling and Human in the Loop
\par}

\vfill

{\large June 2026\par}

\end{center}
\end{titlepage}

\chapter*{Abstract}
\addcontentsline{toc}{chapter}{Abstract}

Thrombosis is a major contributor to cardiovascular morbidity and mortality worldwide. Biomedical knowledge related to thrombosis risk factors evolves continuously through new scientific publications, yet clinicians and researchers primarily rely on static literature reviews and keyword based search systems. These traditional approaches lack semantic understanding, structured retrieval, automated knowledge freshness monitoring, and integrated validation workflows.

This thesis presents the design, implementation, and evaluation of a domain adapted Retrieval Augmented Generation platform tailored for thrombosis risk factor extraction from PubMed literature. The proposed system integrates semantic embedding based retrieval, vector database storage, configurable large language model generation, and a continuous DOI level ingestion scheduler. A human in the loop curation workflow is incorporated through annotation tooling, enabling validation and iterative improvement of extracted knowledge.

The platform follows modular architectural principles emphasizing reproducibility, observability, and configuration driven experimentation. A multi phase evaluation framework assesses retrieval quality, answer grounding, extraction accuracy, and operational reproducibility using qrels based ranking metrics, RAGAS metrics, extraction metrics, and recorded run artifacts. The June 2026 generation benchmark evaluates 25 thrombosis oriented questions across LLM only, RAG without HyDE, and RAG with HyDE modes, with extraction metrics reported only for the 13 questions that have derived gold factors, while the finalized June 2026 retrieval benchmark reports qrels based ranking metrics from a separate 25-query independent retrieval phase.
% evidence: evaluation_metrics/runs/old/june-07-valid-full-run-25-queries/20260605_123942_nogit/paper_benchmark_summary.json

The thesis contributes a reproducible biomedical RAG architecture, an incremental PubMed scheduling mechanism, a human in the loop integration pipeline, and a structured evaluation methodology. The empirical findings are interpreted conservatively: the platform demonstrates auditable experimental architecture and reproducible benchmark artifacts, not clinical readiness.

\clearpage

\chapter*{Acknowledgements}
\addcontentsline{toc}{chapter}{Acknowledgements}

I would like to express my sincere gratitude to my supervisor, Professor Eleni Kaldoudi, for her guidance, academic insight, and continuous support throughout the development of this thesis.

I also thank the supervising PhD candidate, Nikolaos Portoakalidis, for his constructive feedback and technical discussions during the implementation and evaluation phases.

I am grateful to the faculty and colleagues of the MSc Biomedical Informatics program for fostering an interdisciplinary academic environment.

Finally, I would like to thank my family and close friends for their continuous support and encouragement. I would also like to thank my dog for patiently tolerating the postponed walks and for providing moments of balance and perspective throughout this journey.

\clearpage
